% Appendix fragment: no preamble or document environment.
% The parent document defines theorem, proposition, lemma and counterexample.

\section{Binary observations, latent targets, and coherent decision value}

Throughout this appendix, a binary target or observation takes values in
$\{-1,+1\}$ unless a counterexample explicitly uses $0/1$ notation.
The target set and nonnegative weights $w_i$ are fixed before the query.
The loss is weighted Hamming loss
$L(t,a)=\sum_i w_i\mathbf{1}\{t_i\ne a_i\}$, and every coordinatewise
binary action is allowed. Write $W=\sum_i w_i$.

\begin{lemma}[PROVED: binary-observation identity]
\label{lem:binary-observation-r3}
For any coherent joint law of $(T,O)$, put
$m=\mathbb{E}T$ and $c=\mathbb{E}(TO)$. The minimum expected error after
observing $O$ is
\[
 r^O=\frac{1-\max\{|m|,|c|\}}{2}.
\]
Consequently, for target coordinates $T_i$ and candidate observation $O_j$,
\[
 V(j)=\frac12\sum_i w_i\bigl(|c_{ij}|-|m_i|\bigr)_+,
 \qquad m_i=\mathbb{E}T_i,\quad c_{ij}=\mathbb{E}(T_iO_j).
\]
\end{lemma}
\begin{proof}
There are exactly four deterministic maps from a binary observation to a
binary action: $+1,-1,O,-O$. Their correlations with $T$ are respectively
$m,-m,c,-c$, and their errors are one half of one minus that correlation.
Randomization cannot improve a linear expected loss. The best correlation
is therefore $\max\{|m|,|c|\}$. Before querying, the best correlation is
$|m|$. Subtract the two errors and sum coordinatewise, which is legitimate
because Hamming loss is additive and the action space is unrestricted.
\end{proof}

\begin{proposition}[KNOWN/PRIOR ART; PROVED: nonnegative coherent Bayes value]
\label{prop:nonnegative-r3}
Let a fixed decision problem have integrable loss $L(T,a)$ and action set
$\mathcal A$. If the same coherent law is used before and after an
observation $O$, the expected posterior Bayes risk is no larger than the
prior Bayes risk. This applies to any loss, provided the posterior action
minimizes that same loss over the same action set.
\end{proposition}
\begin{proof}
For every prequery action $a\in\mathcal A$, the constant posterior policy
that ignores $O$ is admissible. Hence
\[
 \mathbb{E}\!\left[\inf_{b\in\mathcal A}
          \mathbb{E}[L(T,b)\mid O]\right]
 \le \mathbb{E}\!\left[\mathbb{E}[L(T,a)\mid O]\right]
 =\mathbb{E}L(T,a).
\]
Take the infimum over $a$. For finite binary Hamming decisions, existence
and measurability of minimizers present no difficulty. Equivalently, for
$q_i=\mathbb{P}(T_i=1)$ and $q_i'=\mathbb{P}(T_i=1\mid O)$, the concavity
of $h(q)=\min(q,1-q)$ and the identity $\mathbb{E}q_i'=q_i$ imply
$\mathbb{E}h(q_i')\le h(q_i)$.
\end{proof}

\begin{proposition}[PROVED: Gaussian/probit target-observation moments]
\label{prop:gaussian-probit-r3}
Suppose $F\sim N(\mu,\Sigma)$ is the current working posterior,
$T_i=\operatorname{sign}F_i$, and
$O_j=\operatorname{sign}(F_j+\eta_j)$ with independent
$\eta_j\sim N(0,\tau^2)$, $\tau>0$. For $s_i^2=\Sigma_{ii}>0$,
$g_j^2=\Sigma_{jj}+\tau^2$, and
$\rho_{ij}=\Sigma_{ij}/(s_i g_j)$,
\begin{align*}
 m_i&=2\Phi(\mu_i/s_i)-1,\\
 c_{ij}
 &=4\Phi_2(\mu_i/s_i,\mu_j/g_j;\rho_{ij})
   -2\Phi(\mu_i/s_i)-2\Phi(\mu_j/g_j)+1.
\end{align*}
In particular, for a centered candidate with standard deviation $s$,
\[
 V_{\rm self}
 =\frac1\pi\arcsin\!\left(\frac{s}{\sqrt{s^2+\tau^2}}\right)
 =\frac1\pi\arctan(s/\tau).
\]
\end{proposition}
\begin{proof}
The pair $(F_i,F_j+\eta_j)$ is Gaussian with the displayed variances and
correlation. Its positive-positive orthant probability is the displayed
bivariate normal CDF, by changing both centered standardized signs.
For two binary signs, direct expansion gives
$\mathbb{E}(TO)=4\mathbb{P}(T=1,O=1)-2\mathbb{P}(T=1)
-2\mathbb{P}(O=1)+1$. This proves the general formulas.
At zero means, the Gaussian quadrant identity
$\Phi_2(0,0;\rho)=1/4+\arcsin(\rho)/(2\pi)$ yields
$m=0$ and $c=2\arcsin(\rho)/\pi$. Apply
Lemma~\ref{lem:binary-observation-r3}. The trigonometric equality follows
because $s,\tau>0$. These are exact statements for the current Gaussian
law and the stated probit experiment; they do not make a logistic
posterior or a fresh Laplace refit exact.
\end{proof}

\begin{counterexample}[COUNTEREXAMPLE; PROVED: logistic and probit margin ratios]
\label{ce:logistic-margin-r3}
For both a fixed logistic and a fixed probit observation channel, there
are two independent latent targets for which observed-label margin
uniquely selects a candidate whose coherent Hamming-value ratio tends
to zero.
\end{counterexample}
\begin{proof}
Let $S$ be a fair sign, let $U$ be an independent sign with
$\mathbb{E}U=2\delta>0$, and take
$F_a=\varepsilon S$, $F_b=M U$. Write the channel as
$\mathbb{E}[O\mid F=f]=a(f)$, with either
$a(f)=\tanh(f/2)$ or $a(f)=2\Phi(f/\tau)-1$.
Choose $M>0$ and $0<2\delta<a(M)$. Then
$\mathbb{E}O_a=0$, while $\mathbb{E}O_b=2\delta a(M)>0$,
so observed-label margin selects $a$ uniquely.
The target means are $m_a=0$, $m_b=2\delta$ and the self-moments are
$c_{aa}=a(\varepsilon)$ and $c_{bb}=a(M)$.
Independence gives
$c_{ba}=m_b\mathbb{E}O_a=0$ and
$c_{ab}=m_a\mathbb{E}O_b=0$, so cross-target gains vanish. Thus
\[
 \frac{V(a)}{V(b)}
 =\frac{a(\varepsilon)}{a(M)-2\delta}\longrightarrow0.
\]
For sufficiently small $\varepsilon$, $b$ is the optimal candidate.
\end{proof}

\begin{proposition}[PROVED: centered Gaussian logistic self-value]
If $F=sZ$ with $Z\sim N(0,1)$ and the observation channel is logistic,
then
\[
 V_{\rm self}(s)=\frac12\mathbb{E}\tanh(s|Z|/2)
 \sim \frac{s}{\sqrt{8\pi}}\qquad(s\downarrow0).
\]
The probit channel with $\tau^2=8/\pi$ has the same first-order limit.
\end{proposition}
\begin{proof}
The target mean is zero by symmetry. Its target-observation moment is
$\mathbb{E}[\operatorname{sign}(F)\tanh(F/2)]
=\mathbb{E}\tanh(s|Z|/2)$, proving the equality.
Since $\tanh x\le x$ for $x\ge0$, dominated convergence after division by
$s$ gives limit $\mathbb{E}|Z|/4=1/\sqrt{8\pi}$.
The probit formula follows from
$\arctan(s/\tau)/\pi\sim s/(\pi\tau)$.
\end{proof}

\begin{counterexample}[COUNTEREXAMPLE; PROVED: latent-scale gauge]
\label{ce:scale-gauge-r3}
Let $G$ be any random field with
$\mathbb{P}(G_i=0)=0$ at all finitely many target coordinates, and put
$F^{(c)}=cG$ for $c>0$. All sign laws and samplewise zero sets are
unchanged by $c$. Nevertheless, under any fixed bounded local channel
with response $a$ continuous at zero and $a(0)=0$, the one-query
Hamming values tend to zero as $c\downarrow0$.
\end{counterexample}
\begin{proof}
Positive scaling leaves every sign and zero set unchanged. Conditional
on $G$, the observation mean is $a(cG_j)$, so
\[
 \mathbb{E}[T_iO_j]
 =\mathbb{E}[\operatorname{sign}(G_i)a(cG_j)]
 \longrightarrow0
\]
by bounded dominated convergence. Lemma~\ref{lem:binary-observation-r3}
then gives $V_c(j)\to0$ for every candidate. In a Gaussian example,
means and covariance scale as $c\mu$ and $c^2\Sigma$; standardized means
$\mu_i/s_i$ remain fixed. In particular a centered sign has entropy one
bit for every $c>0$, even though its latent standard deviation vanishes.
When defined, the local boundary-displacement proxy
$s_i/|\nabla\mu_i|$ is also invariant. Under the different, noiseless
channel $O_j=T_j$, all query values are invariant rather than vanishing.
\end{proof}

\begin{theorem}[PROVED: sharp margin guarantee under common sign noise]
\label{thm:bsc-margin-r3}
Suppose the $N\ge2$ candidates are exactly the $N$ unit-weight targets.
For each candidate let $O_j=T_jB_j$, where $B_j$ is independent of the
entire target field and $\mathbb{E}B_j=\alpha\in(0,1]$ is common to all
candidates. Every observed-label margin maximizer $m$ satisfies
\[
 V(m)\ge\frac{1}{N-1}\max_k V(k).
\]
The constant is sharp for every fixed $\alpha>0$.
\end{theorem}
\begin{proof}
Put $a_i=|\mathbb{E}T_i|$ and $q_{ik}=|\mathbb{E}T_iT_k|$.
Since $\mathbb{E}O_j=\alpha\mathbb{E}T_j$, margin minimizes $a_j$.
The value identity becomes
$2V(k)=\sum_i(\alpha q_{ik}-a_i)_+$.
Let $f_b(x)=(b-x)_+$, $s_i=f_\alpha(a_i)$, and
$d_{ik}=f_{\alpha q_{ik}}(a_i)$. For $a_m\le a_k$ and $b\le\alpha$,
\[
 f_\alpha(a_m)-f_\alpha(a_k)
 \ge f_b(a_m)-f_b(a_k).
\]
Indeed the left and right sides are the lengths of the intersections of
$[a_m,a_k]$ with $[0,\alpha]$ and $[0,b]$, respectively.
Using $b=\alpha q_{mk}$ and symmetry of $q_{mk}$ gives
$s_k+d_{mk}\le s_m+d_{km}$.
Every remaining term satisfies $d_{ik}\le s_i\le s_m$, whence
\[
 2V(k)
 \le s_m+d_{km}+(N-2)s_m
 \le (N-1)\,2V(m).
\]
If $V(m)=0$, this inequality also proves all values are zero.

For sharpness when $N\ge3$, take one independent fair decoy sign and
$N-1$ identical signs of mean $2\eta>0$. Margin uniquely selects the
decoy. Its value is $\alpha/2$, whereas a cluster query has value
$(N-1)(\alpha-2\eta)/2$ when $2\eta<\alpha$. For sufficiently small
$\eta$, the cluster query is optimal, and the ratio approaches
$1/(N-1)$ as $\eta\downarrow0$. For $N=2$, the proved constant is one
and cannot be improved. Independence among the different noise variables
$B_j$ is not needed because only one observation is acquired.
\end{proof}

\begin{counterexample}[COUNTEREXAMPLE; PROVED: sign calibration misses amplitude]
\label{ce:amplitude-swap-r3}
For a fixed symmetric logistic or probit channel, identical full
latent-sign laws and identical observation marginals need not yield
identical optimal queries.
\end{counterexample}
\begin{proof}
Take independent centered Gaussian coordinates with standard deviations
$(1,\varepsilon)$ under $P$ and $(\varepsilon,1)$ under $Q$.
Under both laws the sign vector consists of two independent fair signs;
all sign marginals and sign-pair probabilities therefore agree exactly.
Symmetry also makes every observation marginal fair.
The self-value of the unit-scale coordinate is strictly positive,
whereas that of the $\varepsilon$-scale coordinate tends to zero by the
preceding propositions. Independence eliminates cross-target gains.
Consequently the optimal query swaps between $P$ and $Q$.
No score depending only on latent-sign joint probabilities, even when
supplemented with observation marginals, detects this discrepancy.
\end{proof}

\section{Transfer between working and comparison laws}

In this section $P$ and $Q$ concern the same binary target coordinates,
candidate observation actions, fixed weights, and loss. They may differ
in latent distribution or observation channel.

\begin{proposition}[PROVED: moment transfer for Bayes query values]
\label{prop:moment-transfer-r3}
Define
\[
 M=\sum_i w_i|m_i^P-m_i^Q|,
 \qquad C_j=\sum_i w_i|c_{ij}^P-c_{ij}^Q|.
\]
Then
$|V_P(j)-V_Q(j)|\le(M+C_j)/2$.
If $j_P$ and $j_Q$ maximize the corresponding values, then
\[
 V_P(j_P)-V_P(j_Q)
 \le M+\frac{C_{j_P}+C_{j_Q}}{2}
 \le M+\max_j C_j.
\]
\end{proposition}
\begin{proof}
The map $(m,c)\mapsto(|c|-|m|)_+$ is one-Lipschitz in each coordinate,
so apply the triangle inequality to the binary-observation formula.
For the second claim insert
$V_Q(j_P)-V_Q(j_Q)\le0$ between the two $P$-values and apply the
per-candidate bound to each remaining difference.
\end{proof}

\begin{theorem}[PROVED: tight terminal-policy transfer]
\label{thm:terminal-transfer-r3}
Let
\[
 D_j=\sum_i w_i
 \max\{|m_i^P-m_i^Q|,\ |c_{ij}^P-c_{ij}^Q|\}.
\]
Let $j_Q$ minimize terminal Bayes risk under $Q$, and after this query
use a $Q$-optimal decision rule $A^Q_{j_Q}(O_{j_Q})$.
Let $j_P$ minimize terminal Bayes risk under $P$, with terminal risk
$r_P^*(j_P)$. Then the complete query-and-decision policy satisfies
\[
 \mathbb{E}_P L(T,A^Q_{j_Q}(O_{j_Q}))-r_P^*(j_P)
 \le \frac{D_{j_Q}+D_{j_P}}{2}
 \le \max_j D_j.
\]
The universal leading constant one in the last bound cannot be reduced.
The same bound controls query-selection regret when the final decision
after $j_Q$ is instead optimized under $P$.
\end{theorem}
\begin{proof}
Fix a candidate $j$ and a deterministic coordinatewise rule.
Each coordinate rule is one of $+1,-1,O_j,-O_j$, so its correlation
with $T_i$ is one of $\pm m_i,\pm c_{ij}$.
Its error differs between $P$ and $Q$ by at most half the maximum
moment difference appearing in $D_j$. Summing gives
\[
 |\operatorname{risk}_P(A_j)-\operatorname{risk}_Q(A_j)|
 \le D_j/2.
\]
Taking minima also yields $|r_P^*(j)-r_Q^*(j)|\le D_j/2$.
Since $j_Q$ is $Q$-optimal,
\begin{align*}
 \operatorname{risk}_P(A^Q_{j_Q})
 &\le r_Q^*(j_Q)+D_{j_Q}/2\\
 &\le r_Q^*(j_P)+D_{j_Q}/2\\
 &\le r_P^*(j_P)+(D_{j_Q}+D_{j_P})/2.
\end{align*}
Optimizing the final action under $P$ can only decrease the left side.

To show sharpness, one target and one candidate suffice.
Choose $0<\epsilon<1/2$ and $0<\delta$ small.
Let $P$ have moments $(m,c,n)=(\epsilon,-\epsilon,0)$ and let $Q$
have $(m,c,n)=(0,\delta,0)$, where $n=\mathbb{E}O$.
Both are valid laws because
$\mathbb{P}(T=t,O=o)=(1+tm+on+toc)/4$ is nonnegative for these choices.
Under $Q$, the unique optimal rule is $A=O$.
Under $P$ its error is $(1+\epsilon)/2$, whereas optimal error is
$(1-\epsilon)/2$. Regret equals $\epsilon$, while
$D=\epsilon+\delta$. Their ratio tends to one as $\delta\downarrow0$.
\end{proof}

\begin{proposition}[PROVED: latent marginals, conditional structure, and channel]
\label{prop:channel-split-r3}
Assume local channels
$\mathbb{E}_R[O_j\mid F]=a_j^R(F_j)$ for $R\in\{P,Q\}$,
with $|a_j^R|\le1$.
Write $\mu_j^R=\operatorname{Law}_R(F_j)$ and
$h_{ij}^R(f)=\mathbb{E}_R[T_i\mid F_j=f]$.
Using total variation distance
$\operatorname{TV}(\mu,\nu)=\sup_A|\mu(A)-\nu(A)|$, one has
\begin{align*}
 |c_{ij}^P-c_{ij}^Q|
 \le{}&2\operatorname{TV}(\mu_j^P,\mu_j^Q)\\
 &+\mathbb{E}_Q|h_{ij}^P(F_j)-h_{ij}^Q(F_j)|\\
 &+\mathbb{E}_Q|a_j^P(F_j)-a_j^Q(F_j)|.
\end{align*}
The terms separate candidate latent marginal error, latent conditional
structure error, and observation channel error. Target marginal errors
remain explicit in the preceding transfer bounds.
\end{proposition}
\begin{proof}
The local channel assumption gives
$c_{ij}^R=\int h_{ij}^R a_j^R\,d\mu_j^R$.
Add and subtract $\int h_{ij}^P a_j^P\,d\mu_j^Q$ and then
$\int h_{ij}^Q a_j^P\,d\mu_j^Q$.
The first difference is bounded by twice total variation because the
integrand has absolute value at most one. The other two differences are
bounded by the displayed expectations, since $|h|,|a|\le1$.
One may choose any bounded version of $h^P$ on $Q$-only support;
the identity and bound remain valid, although the bound can then be
uninformative. The conditional-structure term is not asserted to be a
pure copula discrepancy.
\end{proof}

\begin{proposition}[PROVED: alternative marginal/dependence/channel split]
Let $\gamma_{ij}^R=\operatorname{Law}_R(F_i,F_j)$ and
$\kappa_{ij}^R=\gamma_{ij}^R-\mu_i^R\otimes\mu_j^R$.
Set $\delta_i=\operatorname{TV}(\mu_i^P,\mu_i^Q)$ and let
$\|\cdot\|_{\rm var}$ denote the full variation norm of a signed measure.
Under the preceding local-channel assumptions,
\begin{align*}
 |c_{ij}^P-c_{ij}^Q|
 \le{}&2\delta_i+2\delta_j+
       \|\kappa_{ij}^P-\kappa_{ij}^Q\|_{\rm var}\\
 &+\mathbb{E}_Q|a_j^P(F_j)-a_j^Q(F_j)|,
 \qquad |m_i^P-m_i^Q|\le2\delta_i.
\end{align*}
\end{proposition}
\begin{proof}
First keep the $P$ channel fixed. The integrand
$\operatorname{sign}(f_i)a_j^P(f_j)$ is bounded by one.
Decompose the pair-law difference into its product-marginal difference
and $\kappa_{ij}^P-\kappa_{ij}^Q$.
The product measures have total variation distance at most
$\delta_i+\delta_j$, as follows by adding and subtracting
$\mu_i^Q\otimes\mu_j^P$. Integration against their difference is
therefore bounded by $2\delta_i+2\delta_j$.
The dependence part is bounded by its full variation norm.
Finally change the channel under $Q$ and use its displayed expectation.
The marginal sign bound follows directly from the bounded sign function.
\end{proof}

\begin{proposition}[PROVED: algorithmic gain versus Bayes information]
\label{prop:updater-decomposition-r3}
For a fixed prequery action $a_0$ and algorithmic postquery rule
$A_j(O_j)$, define
\[
 b_P=\operatorname{risk}_P(a_0)-r_P^0,\qquad
 e_P(j)=\operatorname{risk}_P(A_j)-r_P^*(j).
\]
Both are nonnegative, and the algorithmic gain is exactly
\[
 G_P^A(j)=V_P(j)+b_P-e_P(j).
\]
Consequently, for any model-selected query $j_Q$,
\[
 \max_jG_P^A(j)-G_P^A(j_Q)
 =\max_j\{V_P(j)-e_P(j)\}
    -\{V_P(j_Q)-e_P(j_Q)\}.
\]
This residual does not, without additional controls, identify
information absent from the working model.
\end{proposition}
\begin{proof}
By definition,
$\operatorname{risk}_P(a_0)=r_P^0+b_P$ and
$\operatorname{risk}_P(A_j)=r_P^*(j)+e_P(j)$.
Subtract and use $V_P(j)=r_P^0-r_P^*(j)$.
The second identity follows because $b_P$ is independent of $j$.
Even if the working law equals $P$, different updater excess risks
$e_P(j)$ can leave positive algorithmic regret.
If $P$ is conditioned on a fully known realized target field, then
$r_P^0=r_P^*(j)=V_P(j)=0$; nevertheless a restricted learner may have
positive gains because its own errors change. Such fixed-truth gains
are algorithmic utilities, not positive Bayes information under the
omniscient conditional law.
\end{proof}

\section{Coherent negative values for other risk functionals}

The next examples use $T_1,T_2\in\{0,1\}$.
A plug-in edge prediction is the exclusive-or of the two marginal
mode predictions. Its edge risk is the probability that this predicted
cut differs from the true cut. Its Dice ratio is
\[
 R_D=\frac{\mathbb{E}W_{\rm err}}
           {\mathbb{E}W_{\rm true}+W_{\rm pred}}.
\]
Neither operation is automatically the minimum posterior expectation
of one fixed loss, so Proposition~\ref{prop:nonnegative-r3} does not
automatically apply.

\begin{counterexample}[COUNTEREXAMPLE; PROVED: negative coherent plug-in edge gain]
Assign masses $(2,2,2,3)/9$ to states $(00,01,10,11)$ and observe
$O=\mathbf{1}\{(T_1,T_2)=(0,1)\}$. Exact Bayesian conditioning gives
an expected plug-in edge-risk reduction of $-1/9$.
\end{counterexample}
\begin{proof}
Before observing, both marginal modes are one, so the edge is predicted
uncut. Its error is the true-cut probability $4/9$.
The branch $O=1$ identifies state $01$ exactly and has zero error.
On $O=0$, whose probability is $7/9$, the conditional masses on
$(00,10,11)$ are $(2,2,3)/7$. The marginal probabilities of one are
$5/7$ and $3/7$, so the plug-in modes are $(1,0)$ and predict a cut.
Its conditional error is the no-cut probability $5/7$.
Expected future error is $(7/9)(5/7)=5/9$, so the reduction is
$4/9-5/9=-1/9$. All marginal decisions are strict, and every update
is exact conditioning.
\end{proof}

\begin{counterexample}[COUNTEREXAMPLE; PROVED: negative coherent Dice-ratio gain]
Assign masses $(1,1,2,1)/5$ to $(00,01,10,11)$ and observe the same
$O=\mathbf{1}\{(T_1,T_2)=(0,1)\}$. Exact conditioning gives an expected
Dice-ratio reduction of $-1/60$.
\end{counterexample}
\begin{proof}
The prior marginal modes are $(1,0)$, the true-cut probability is $3/5$,
and the ratio is
$R_D=(2/5)/(1+3/5)=1/4$.
The $O=1$ branch has zero error and zero ratio.
On $O=0$, of probability $4/5$, conditional masses on $(00,10,11)$
are $(1,2,1)/4$. Marginal modes remain strictly $(1,0)$.
The true-cut probability is $1/2$, so
$R_D=(1/2)/(1+1/2)=1/3$.
Expected future ratio is $(4/5)(1/3)=4/15$ and the reduction is
$1/4-4/15=-1/60$.
Thus a negative ratio reduction by itself does not demonstrate an
incoherent posterior. A separately measured failure of the posterior
martingale identity is a different diagnostic.
\end{proof}
